% !TeX root = ../main.tex
\chapter{Hướng dẫn chạy chương trình}
\label{chap:instructions}

\section{Yêu cầu môi trường}

Notebook sử dụng Python 3 cùng các package trong Bảng~\ref{tab:environment}.
Dataset được scikit-learn cung cấp sẵn nên không cần tải file thủ công.

\begin{table}[H]
\centering
\caption{Môi trường thực thi notebook}
\label{tab:environment}
\begin{tabular}{ll}
\toprule
Thành phần & Vai trò \\
\midrule
Python 3.10 trở lên & Runtime \\
NumPy & Tính toán số và danh sách alpha \\
pandas & Bảng dữ liệu và xuất CSV \\
Matplotlib & Trực quan hóa \\
scikit-learn & Dataset, Decision Tree, GridSearchCV và metrics \\
Jupyter / VS Code & Mở và chạy notebook \\
\bottomrule
\end{tabular}
\end{table}

\section{Các bước tái lập}

Tại thư mục chứa notebook, tạo môi trường và cài package:

\begin{verbatim}
python -m venv .venv
.\.venv\Scripts\Activate.ps1
python -m pip install -r requirements.txt
\end{verbatim}

Mở \texttt{Decision\_Tree\_2.ipynb}, chọn kernel của môi trường vừa tạo, sau
đó chạy \textbf{Restart Kernel and Run All Cells}. Notebook phải chạy theo thứ
tự từ trên xuống vì các model cải thiện tái sử dụng cùng training/test split.

\section{Đầu ra kỹ thuật}

Sau khi chạy thành công, notebook tạo:

\begin{itemize}
  \item \texttt{results/model\_comparison.csv}: toàn bộ metric và complexity;
  \item \texttt{results/hyperparameter\_summary.csv}: tham số được chọn và CV
  accuracy;
  \item \texttt{results/SUMMARY.md}: dataset, protocol, kết quả khóa và winner;
  \item \texttt{results/figures/}: confusion matrix, full/readable baseline,
  feature importance, tuning curves, bảng so sánh dạng hình và ba cây cải thiện.
\end{itemize}

Không sửa số liệu thủ công trong CSV hoặc hình. Khi seed, split hoặc parameter
grid thay đổi, cần chạy lại toàn bộ notebook và cập nhật đồng thời report/video.

\section{Kiểm tra trước khi bàn giao}

Notebook đã được thực thi từ đầu đến cuối và lưu kèm output. Người nhận nên xác
nhận các giá trị chính: baseline accuracy 0,9123; Max Depth và Pruning accuracy
0,9386; selected model là Max Depth. File report dùng đường dẫn tương đối đến
\texttt{../results/figures}, nên giữ nguyên cấu trúc thư mục khi build LaTeX.
